\section{Results}

% ---------- Slide 15 ----------
\begin{frame}{To hide, the adversary mostly cuts itself off}
  \SectionPill{Result 1: decrease}{PresFourDecrease}
  \begin{columns}[T,onlytextwidth]
    \begin{column}{0.50\textwidth}
      \centering
      \begin{tikzpicture}
        \node[modelbox,minimum width=4.2cm,minimum height=1.7cm,inner sep=3pt] (m1) at (0,1.0) {
          \begin{tikzpicture}[scale=0.45]
            \node[gadv] (a) at (0,0) {};
            \node[gnew] (n) at (1.4,0) {};
            \draw[cutedge] (a)--(n);
            \node[gnode] (f1) at (-1.1,0.8) {};
            \node[gnode] (f2) at (-1.1,-0.8) {};
            \draw[gedge] (a)--(f1); \draw[gedge] (a)--(f2);
          \end{tikzpicture}\\[-0.15em]
          {\scriptsize\textbf{remove\_node\_to\_self}}
        };
        \node[modelbox,minimum width=4.2cm,minimum height=1.7cm,inner sep=3pt] (m2) at (0,-1.1) {
          \begin{tikzpicture}[scale=0.45]
            \node[gadv] (a) at (0,0) {};
            \node[gnode] (f) at (1.4,0) {};
            \draw[cutedge] (a)--(f);
            \node[gnode] (f1) at (-1.1,0.8) {};
            \node[gnode] (f2) at (-1.1,-0.8) {};
            \draw[gedge] (a)--(f1); \draw[gedge] (a)--(f2);
          \end{tikzpicture}\\[-0.15em]
          {\scriptsize\textbf{self\_remove\_edge\_friend}}
        };
      \end{tikzpicture}
    \end{column}
    \begin{column}{0.48\textwidth}
      \centering
      \includegraphics[width=\linewidth]{fig3_random.png}\\[0.10em]
      {\scriptsize\color{PresFourMuted}Fig. 3 --- random networks. \textcolor{PresFourDecrease}{Blue} / \textcolor{PresFourDecreaseLocal}{teal} bars are decrease attacks; note the long bars on \texttt{self\_remove\_edge\_friend}.}
    \end{column}
  \end{columns}
  \vspace{0.05em}
  \begin{beamercolorbox}[wd=\linewidth,rounded=true,sep=0.40em]{takeaway box}
    {\color{PresFourDecrease}\bfseries Hiding has a simple recipe}\quad
    Drop the most informative tie or neighbor; same two moves dominate across topologies and centralities. Local moves are often enough.
  \end{beamercolorbox}
  \SlideNote{3:00}{Main finding: decreasing centrality uses a smaller, more consistent set of moves than increasing.}
\end{frame}

% ---------- Slide 16 ----------
\begin{frame}{To look more powerful, the best move depends on context}
  \SectionPill{Result 2: increase}{PresFourIncrease}
  \begin{columns}[T,onlytextwidth]
    \begin{column}{0.55\textwidth}
      \centering
      \includegraphics[width=\linewidth]{fig6_cellular.png}\\[0.10em]
      {\scriptsize\color{PresFourMuted}Fig. 6 --- cellular networks. \textcolor{PresFourIncrease}{Red} / \textcolor{PresFourIncreaseLocal}{pink} bars spread across many moves; the best one shifts per centrality.}
    \end{column}
    \begin{column}{0.43\textwidth}
      \MiniCard{PresFourIncrease}{add\_edge\_anywhere}{Reach beyond the ego network.}
      \vspace{0.15em}
      \MiniCard{PresFourIncrease}{self\_add\_edge\_foaf}{Connect to a friend-of-friend.}
      \vspace{0.15em}
      \MiniCard{PresFourIncrease}{add\_node\_to\_self}{Attach a brand-new node.}
    \end{column}
  \end{columns}
  \vspace{0.10em}
  {\small\color{PresFourMuted}Boosting is varied; hiding is concentrated. Global reach matters more for boosting.}
  \SlideNote{2:30}{Hiding has a simple strategy; boosting is context-dependent.}
\end{frame}

% ---------- Slide 17 ----------
\begin{frame}{Global moves matter more for boosting than hiding}
  \SectionPill{Result 3: magnitude}{PresFourAccent}
  \begin{columns}[T,onlytextwidth]
    \begin{column}{0.66\textwidth}
      \centering
      \includegraphics[width=\linewidth]{fig7_distribution.png}\\[0.10em]
      {\scriptsize\color{PresFourMuted}Fig. 7 --- distribution of centrality change. Left of 0 = decrease, right of 0 = increase.}
    \end{column}
    \begin{column}{0.32\textwidth}
      \MiniCard{PresFourDecrease}{Decrease}{Local and global overlap --- one cut is often enough.}
      \vspace{0.15em}
      \MiniCard{PresFourIncrease}{Increase}{Global moves create \emph{larger} gains for closeness.}
      \vspace{0.15em}
      \MiniCard{PresFourAccent}{Degree omitted}{Changes are $\pm 1$ --- trivial distribution.}
    \end{column}
  \end{columns}
  \SlideNote{2:30}{Move from `which move was optimal?' to `how large was the effect?'. Figure 7 shows the score distribution.}
\end{frame}
